% ============================================================================
% CHAPTER 3: CONTRACT FRAMEWORK
% ============================================================================

\chapter{The Contract Framework}
\label{ch:contracts}

\importantbox{
\textbf{Contracts are the cornerstone of this system's safety guarantees.} Every component operates under formally specified Assume-Guarantee (A/G) contracts that enable compositional verification and principled adaptive behavior.
}

\section{Foundations of Assume-Guarantee Contracts}

\subsection{Formal Definition}

\begin{definition}[Assume-Guarantee Contract]
\label{def:ag_contract}
An Assume-Guarantee (A/G) contract is a tuple $\mathcal{C} = (\mathcal{V}, \mathcal{A}, \mathcal{G})$ where:
\begin{itemize}
    \item $\mathcal{V}$ is a set of variables (inputs and outputs)
    \item $\mathcal{A}$ is a set of \textbf{assumptions} over input variables
    \item $\mathcal{G}$ is a set of \textbf{guarantees} over output variables
\end{itemize}
\end{definition}

The semantics of a contract are captured by the implication:
\begin{equation}
\boxed{\mathcal{A} \implies \mathcal{G}}
\label{eq:contract_semantics}
\end{equation}

\begin{theorem}[Contract Satisfaction]
A component $M$ \textbf{satisfies} contract $\mathcal{C} = (\mathcal{A}, \mathcal{G})$, written $M \models \mathcal{C}$, if and only if:
\begin{equation}
\forall \text{ inputs satisfying } \mathcal{A}: \text{ outputs of } M \text{ satisfy } \mathcal{G}
\end{equation}
\end{theorem}

\subsection{Why Contracts Matter for Safety}

Traditional safety approaches suffer from several limitations:

\begin{table}[htbp]
    \centering
    \caption{Comparison of safety approaches}
    \label{tab:safety_comparison}
    \begin{tabular}{@{}lcc@{}}
        \toprule
        \textbf{Aspect} & \textbf{Heuristic Safety} & \textbf{Contract-Based} \\
        \midrule
        Formal verification & \textcolor{red}{\ding{55}} No & \textcolor{green}{\ding{51}} Yes \\
        Compositional & \textcolor{red}{\ding{55}} No & \textcolor{green}{\ding{51}} Yes \\
        Transparent reasoning & \textcolor{red}{\ding{55}} Limited & \textcolor{green}{\ding{51}} Full \\
        Graceful degradation & \textcolor{red}{\ding{55}} Ad-hoc & \textcolor{green}{\ding{51}} Principled \\
        Pre-flight verification & \textcolor{red}{\ding{55}} No & \textcolor{green}{\ding{51}} Yes \\
        \bottomrule
    \end{tabular}
\end{table}

\section{Contract Composition with Pacti}

\subsection{The Pacti Library}

We use the \textbf{Pacti library}\footnote{\url{https://pacti.readthedocs.io/}} for polyhedral contract manipulation. Pacti represents contracts as systems of linear inequalities, enabling:

\begin{itemize}
    \item \textbf{Contract composition}: Combining component contracts
    \item \textbf{Contract quotient}: Deriving required component contracts
    \item \textbf{Contract merging}: Combining multiple viewpoints
    \item \textbf{Refinement checking}: Verifying contract relationships
\end{itemize}

\subsection{Polyhedral Contracts}

\begin{definition}[Polyhedral Contract]
A polyhedral contract represents assumptions and guarantees as polyhedra defined by linear inequalities:
\begin{align}
    \mathcal{A} &: \mathbf{A}_a \mathbf{v} \leq \mathbf{b}_a \\
    \mathcal{G} &: \mathbf{A}_g \mathbf{v} \leq \mathbf{b}_g
\end{align}
where $\mathbf{v}$ is the vector of contract variables.
\end{definition}

\begin{example}[PID Controller Contract in Polyhedral Form]
The PID controller contract can be expressed as:
\begin{align}
    \mathcal{A}_{\text{PID}}: \quad & \text{wind\_speed} \leq 3.0 \\
    & \text{position\_error} \leq 2.0 \\
    & \text{disturbance} \leq 1.0 \\[6pt]
    \mathcal{G}_{\text{PID}}: \quad & \text{settling\_time} \leq 5.0 \\
    & \text{overshoot} \leq 0.3 \\
    & \text{steady\_state\_error} \leq 0.5
\end{align}
\end{example}

\subsection{Contract Composition}

When two components are connected in series, their contracts compose:

\begin{theorem}[Contract Composition]
\label{thm:composition}
Given contracts $\mathcal{C}_1 = (\mathcal{A}_1, \mathcal{G}_1)$ and $\mathcal{C}_2 = (\mathcal{A}_2, \mathcal{G}_2)$ where outputs of $\mathcal{C}_1$ feed inputs of $\mathcal{C}_2$, the composed contract $\mathcal{C}_1 \otimes \mathcal{C}_2$ is:
\begin{equation}
\mathcal{C}_1 \otimes \mathcal{C}_2 = \left( \mathcal{A}_1 \cup (\mathcal{A}_2 \setminus \mathcal{G}_1), \; \mathcal{G}_1 \cap \mathcal{G}_2 \right)
\label{eq:composition}
\end{equation}
\end{theorem}

\begin{figure}[htbp]
    \centering
    \begin{tikzpicture}[
        contract/.style={
            draw,
            rounded corners,
            minimum width=3cm,
            minimum height=2cm,
            align=center
        }
    ]
        % Contract 1
        \node[contract, fill=pidblue!20] (c1) at (0,0) {
            $\mathcal{C}_1$\\[3pt]
            $(\mathcal{A}_1, \mathcal{G}_1)$
        };

        % Contract 2
        \node[contract, fill=hinfgreen!20] (c2) at (5,0) {
            $\mathcal{C}_2$\\[3pt]
            $(\mathcal{A}_2, \mathcal{G}_2)$
        };

        % Composed contract
        \node[contract, fill=contractpurple!20, minimum width=4cm] (comp) at (2.5,-3.5) {
            $\mathcal{C}_1 \otimes \mathcal{C}_2$\\[3pt]
            $(\mathcal{A}_1 \cup (\mathcal{A}_2 \setminus \mathcal{G}_1), \mathcal{G}_1 \cap \mathcal{G}_2)$
        };

        % Arrows
        \draw[-{Stealth}, thick] (c1) -- (c2) node[midway, above] {$\mathcal{G}_1 \to \mathcal{A}_2$};
        \draw[-{Stealth}, thick, dashed] (c1.south) -- (comp.north west);
        \draw[-{Stealth}, thick, dashed] (c2.south) -- (comp.north east);
        \node at (2.5,-1.5) {compose};
    \end{tikzpicture}
    \caption{Contract composition: outputs of $\mathcal{C}_1$ satisfy inputs of $\mathcal{C}_2$}
    \label{fig:composition}
\end{figure}

\section{System Contract Hierarchy}

\subsection{Sensor Contracts}

\begin{contract}[GPS Sensor Contract]
\label{contract:gps}
\begin{align*}
    \mathcal{A}_{\text{GPS}}: \quad & \text{satellites} \geq 6 \\
    & \text{hdop} \leq 2.0 \\
    & \text{sky\_visibility} > 0.5 \\[6pt]
    \mathcal{G}_{\text{GPS}}: \quad & \text{position\_error} \leq 2.0 \text{ m} \\
    & \text{velocity\_error} \leq 0.5 \text{ m/s} \\
    & \text{update\_rate} \geq 10 \text{ Hz}
\end{align*}
\end{contract}

\begin{contract}[IMU Sensor Contract]
\label{contract:imu}
\begin{align*}
    \mathcal{A}_{\text{IMU}}: \quad & \text{calibrated} = \text{true} \\
    & |\text{temperature} - 25| \leq 10 \text{ °C} \\[6pt]
    \mathcal{G}_{\text{IMU}}: \quad & \text{attitude\_error} \leq 0.1 \text{ rad} \\
    & \text{rate\_error} \leq 0.05 \text{ rad/s} \\
    & \text{update\_rate} \geq 200 \text{ Hz}
\end{align*}
\end{contract}

\subsection{Estimator Contract}

\begin{contract}[EKF Estimator Contract]
\label{contract:ekf}
\begin{align*}
    \mathcal{A}_{\text{EKF}}: \quad & \mathcal{G}_{\text{GPS}} \lor \mathcal{G}_{\text{IMU}} \text{ (at least one sensor)} \\
    & \text{computation\_time} \leq 10 \text{ ms} \\[6pt]
    \mathcal{G}_{\text{EKF}}: \quad & \text{state\_estimate\_valid} = \text{true} \\
    & \text{covariance\_bounded} = \text{true} \\
    & \text{position\_uncertainty} \leq 5.0 \text{ m (degraded mode)}
\end{align*}
\end{contract}

\subsection{Controller Contracts}

\begin{contract}[PID Controller Contract]
\label{contract:pid}
The PID controller operates efficiently under nominal conditions:
\begin{align*}
    \mathcal{A}_{\text{PID}}: \quad & \text{wind\_speed} \leq 3.0 \text{ m/s} \\
    & \text{position\_error} \leq 2.0 \text{ m} \\
    & \text{velocity\_error} \leq 1.0 \text{ m/s} \\
    & \text{disturbance} \leq 1.0 \\
    & \text{state\_estimate\_valid} = \text{true} \\[6pt]
    \mathcal{G}_{\text{PID}}: \quad & \text{settling\_time} \leq 5.0 \text{ s} \\
    & \text{overshoot} \leq 30\% \\
    & \text{steady\_state\_error} \leq 0.5 \text{ m} \\
    & \text{control\_effort} \leq \text{nominal}
\end{align*}
\end{contract}

\begin{contract}[H-infinity Controller Contract]
\label{contract:hinf}
The H-infinity controller provides robust stabilization with minimal assumptions:
\begin{align*}
    \mathcal{A}_{\text{H}\infty}: \quad & \text{wind\_speed} \leq 8.0 \text{ m/s} \\
    & \text{(no position error assumption)} \\
    & \text{(no disturbance assumption)} \\[6pt]
    \mathcal{G}_{\text{H}\infty}: \quad & \text{settling\_time} \leq 10.0 \text{ s} \\
    & \text{stability} = \text{guaranteed} \\
    & \text{altitude\_maintained} = \text{true} \\
    & \text{attitude\_leveled} = \text{true}
\end{align*}
\end{contract}

\importantbox{
\textbf{Key Observation:} The H-infinity contract has \textit{weaker assumptions} (accepts higher wind, no position error requirement) but provides \textit{weaker guarantees} (slower settling, no overshoot bound). This tradeoff enables principled controller switching.
}

\subsection{Safety Contract}

\begin{contract}[CBF Safety Contract]
\label{contract:cbf}
The CBF filter enforces hard constraints with \textbf{no assumptions}:
\begin{align*}
    \mathcal{A}_{\text{CBF}}: \quad & \text{(none --- always active)} \\[6pt]
    \mathcal{G}_{\text{CBF}}: \quad & z \in [-50, -2] \text{ m (altitude bounds)} \\
    & \|\mathbf{v}\| \leq 15 \text{ m/s (velocity limit)} \\
    & |\phi|, |\theta| \leq 30° \text{ (tilt limit)} \\
    & \|[\,p, q, r\,]\| \leq 3 \text{ rad/s (rate limit)}
\end{align*}
\end{contract}

\section{Contract-Based Controller Selection}

The contract framework enables \textbf{principled} controller selection based on assumption satisfaction:

\begin{algorithm}[htbp]
\caption{Contract-Based Controller Selection}
\label{alg:selection}
\begin{algorithmic}[1]
    \Require Current system conditions $\mathbf{c}$
    \Ensure Selected controller name
    \State $\text{pid\_assumptions} \gets$ CheckAssumptions($\mathcal{C}_{\text{PID}}$, $\mathbf{c}$)
    \State $\text{hinf\_assumptions} \gets$ CheckAssumptions($\mathcal{C}_{\text{H}\infty}$, $\mathbf{c}$)
    \If{$\text{pid\_assumptions} = \text{satisfied}$}
        \State \Return ``PID'' \Comment{Use efficient controller}
    \ElsIf{$\text{hinf\_assumptions} = \text{satisfied}$}
        \State \Return ``H-infinity'' \Comment{Use robust controller}
    \Else
        \State \Return ``H-infinity'' + TriggerEmergency()
    \EndIf
\end{algorithmic}
\end{algorithm}

\begin{figure}[htbp]
    \centering
    \begin{tikzpicture}[
        node distance=1.5cm,
        decision/.style={
            diamond,
            draw,
            fill=contractpurple!20,
            text width=2.8cm,
            align=center,
            inner sep=2pt
        },
        block/.style={
            rectangle,
            draw,
            rounded corners,
            fill=white,
            minimum width=2.5cm,
            minimum height=1cm,
            align=center
        },
        arrow/.style={-{Stealth}, thick}
    ]
        % Nodes
        \node[block, fill=gray!20] (start) {System\\Conditions};
        \node[decision, below=of start] (pid_check) {PID\\assumptions\\satisfied?};
        \node[block, fill=pidblue!30, below left=1.5cm and 1.5cm of pid_check] (pid) {Use PID\\(Efficient)};
        \node[decision, below right=1.5cm and 1.5cm of pid_check] (hinf_check) {H-$\infty$\\assumptions\\satisfied?};
        \node[block, fill=hinfgreen!30, below left=1.5cm and 0cm of hinf_check] (hinf) {Use H-$\infty$\\(Robust)};
        \node[block, fill=safetyred!30, below right=1.5cm and 0cm of hinf_check] (emergency) {Emergency\\Mode};

        % Arrows
        \draw[arrow] (start) -- (pid_check);
        \draw[arrow] (pid_check) -- node[left] {Yes} (pid);
        \draw[arrow] (pid_check) -- node[right] {No} (hinf_check);
        \draw[arrow] (hinf_check) -- node[left] {Yes} (hinf);
        \draw[arrow] (hinf_check) -- node[right] {No} (emergency);
    \end{tikzpicture}
    \caption{Contract-based controller selection decision tree}
    \label{fig:selection}
\end{figure}

\section{Runtime Contract Monitoring}

During flight, contracts are continuously monitored:

\begin{lstlisting}[caption={Runtime contract monitoring}]
def monitor_contracts(self, conditions, timestamp):
    """Monitor all contracts and log violations"""

    # Check PID assumptions
    pid_contract = self.contracts['PID']
    pid_satisfied, pid_msg = pid_contract.check_assumptions(conditions)

    # Check H-infinity assumptions
    hinf_contract = self.contracts['Hinf']
    hinf_satisfied, hinf_msg = hinf_contract.check_assumptions(conditions)

    # Log status
    metrics = ContractMetrics(
        timestamp=timestamp,
        assumptions_met=pid_satisfied or hinf_satisfied,
        active_contract='PID' if pid_satisfied else 'Hinf'
    )

    return metrics
\end{lstlisting}

\section{Pre-flight Contract Verification}

Before takeoff, the system composes contracts to verify mission feasibility:

\begin{theorem}[Mission Feasibility]
A mission $\mathcal{M}$ with initial conditions $\mathbf{c}_0$ is feasible if:
\begin{equation}
\exists \mathcal{C} \in \{\mathcal{C}_{\text{PID}}, \mathcal{C}_{\text{H}\infty}\}: \mathbf{c}_0 \models \mathcal{A}_{\mathcal{C}}
\end{equation}
That is, at least one controller's assumptions are satisfied by initial conditions.
\end{theorem}

\begin{lstlisting}[caption={Pre-flight contract verification}]
def pre_flight_check(self, conditions, mission):
    """Verify mission feasibility before takeoff"""

    # Try PID first (most efficient)
    feasible_pid = self.verify_mission_feasibility('PID', conditions)
    if feasible_pid:
        return True, "Mission feasible with PID"

    # Try H-infinity (always works if wind < 8 m/s)
    feasible_hinf = self.verify_mission_feasibility('Hinf', conditions)
    if feasible_hinf:
        return True, "Mission feasible with H-infinity"

    # Neither controller can handle conditions
    return False, "Mission infeasible - conditions too extreme"
\end{lstlisting}
